%% ===================================================================
%% PART II — Logbook addendum: pivot to unified generative MM (Janus)
%% ADDITIVE ONLY. Paste AFTER Ivan's existing sections, BEFORE the
%% bibliography. Do not modify or delete anything already written.
%% Style matches the existing document (booktabs, multirow, natbib).
%% Each section below is a logbook entry: coherent, validated, concise.
%% ===================================================================

%% -------------------------------------------------------------------
\section{Project Evolution: From Dense-Prediction 4M to Unified Generative Models}
%% -------------------------------------------------------------------

Part~I studied PEFT on 4M-21, a native multimodal model whose modalities are
image-derived dense maps (depth, surface normals, semantic segmentation) plus
captioning. Over the course of those experiments the research question itself
sharpened: the interaction we want to characterise is text$\leftrightarrow$image
\emph{understanding and generation} within a single shared transformer ---
the axis along which open native multimodal models are currently most active,
and the one in which the unified-model thesis (one backbone, many modalities)
is genuinely tested. 4M-21 does not exercise this axis: it has no native
text-to-image generation. Part~II is therefore a logbook of the pivot toward
\emph{natively unified} understanding+generation models, of the candidates we
considered, and of the pipeline we built and validated on the chosen base
model. We report what worked, what did not, and the first preliminary results;
findings from Part~I are preserved as the methodological foundation on which
the new experiments rest.

%% -------------------------------------------------------------------
\section{Model Selection: A Design-Space Exploration}
%% -------------------------------------------------------------------

After the 4M-21 phase we re-examined which base model best fits a research
question centred on the understanding$\leftrightarrow$generation interaction.
Below we state the selection criteria and the reasons each candidate was kept
or set aside. The discarded models are themselves a useful finding: they map
out which open native multimodal models are, in practice, amenable to
gradient-based PEFT today.

\subsection{Selection criteria}
A base model suitable for our experiments must:
\begin{enumerate}
    \item be \emph{natively unified} --- perform both multimodal understanding
          (image$\rightarrow$text) and generation (text$\rightarrow$image)
          within a single shared transformer backbone;
    \item expose its weights as a standard \texttt{nn.Module}, so that
          gradient-based PEFT (e.g.\ \texttt{peft.get\_peft\_model}) applies
          directly without re-implementing the model;
    \item have image generation \emph{actually enabled} in the open release;
    \item be openly available and not gated;
    \item fit our compute budget ($4\times$A100) at a tractable size.
\end{enumerate}

\subsection{4M-21: exploratory phase and lessons}
The 4M-21 phase (Part~I) established methodology we still rely on: LoRA
injection on attention projections and a cross-task transfer protocol for
measuring forgetting. It also surfaced lessons that shaped the present work ---
a cross-entropy-vs-metric disconnect on depth (token-level CE decreasing
without a matching gain in metric-space depth accuracy), a train/validation
generalisation gap diagnosed via an overfit test, and two pipeline bugs (a
no-gradient crash on all-padding batches and an inconsistent train/validation
loss masking) whose fixes we carry forward. The decisive limitation, however,
is conceptual rather than technical: 4M-21's modalities are image-derived dense
maps (depth, surface normals, semantic segmentation), not natural-language
$\leftrightarrow$ image generation. It therefore does not exercise the
understanding$\leftrightarrow$generation tension that defines our question.
% TODO (Francesco to provide numbers): overfit test (CE-->0 vs generalisation),
% AbsRel/delta on depth, normals forgetting on DIODE. Quantitative 4M results
% live in Part~I (Table~\ref{tab:results}); extra figures to be added here.

\subsection{Chameleon: image generation disabled in the open release}
Chameleon \citep{chameleon2024} is an early-fusion, fully-tokenised unified
model --- the right family. However, image generation is disabled in the public
Hugging Face release: the detokenisation path that turns generated image tokens
back into pixels is not exposed, so the generation direction cannot be studied
without modifying the model. Criterion~(iii) fails.

\subsection{Anole: inference-only architecture blocks PEFT}
Anole \citep{chern2024anole} re-enables generation on Chameleon by selectively
unfreezing the image-token output logits --- conceptually close to a
modality-aware PEFT. Its open implementation, however, executes the model inside
a spawned worker subprocess for inference; the parameters are not reachable as a
live \texttt{nn.Module} in the main process. Gradient-based PEFT is therefore
not possible without re-implementing the serving stack. Criterion~(ii) fails.

\subsection{Janus-Pro: the tractable choice}
Janus-Pro \citep{wu2024janus} satisfies every criterion. It is integrated in
Hugging Face \texttt{transformers} as \texttt{JanusForConditionalGeneration} (a
standard \texttt{nn.Module}); it performs both understanding
(\texttt{generation\_mode="text"}) and generation
(\texttt{generation\_mode="image"}); it \emph{decouples} the visual encoders ---
a SigLIP encoder for understanding and a VQ-VAE for generation --- over a single
shared transformer; it is released under MIT/DeepSeek licences and is not gated;
and it is available at 1B and 7B parameters. We prototype on the 1B model and
reserve the 7B for final numbers. The decoupled encoders also make Janus the
natural testbed for our LoRA-placement question (Section~\ref{sec:janus_rqs}).

\begin{table}[h]
\centering
\small
\begin{tabular}{lccc}
\toprule
Requirement & Chameleon & Anole & Janus-Pro \\
\midrule
Native HF \texttt{nn.Module}        & Yes & No (subprocess) & Yes \\
Weights trainable for PEFT          & Yes & No              & Yes \\
Image generation enabled (open)     & No  & Yes             & Yes \\
Decoupled understanding/gen.\ enc.  & No  & No              & Yes \\
Open / not gated                    & Yes & Yes             & Yes \\
\bottomrule
\end{tabular}
\caption{Why each open native multimodal model was kept or ruled out for
gradient-based PEFT experiments. Janus-Pro is the only candidate that satisfies
all criteria.}
\label{tab:model_selection}
\end{table}

%% -------------------------------------------------------------------
\section{Janus-Pro: Experimental Setup}
\label{sec:janus_setup}
%% -------------------------------------------------------------------

\paragraph{Model.} We use \texttt{deepseek-community/Janus-Pro-1B} for
prototyping and reserve the 7B variant for the final, scaling-comparable
numbers. Both share the decoupled-encoder architecture: a SigLIP visual encoder
for understanding, a VQ-VAE (codebook size $16{,}384$, $576$ tokens per image)
for generation, over a single shared Llama-style transformer backbone.

\paragraph{Data.} We use COCO val2017 (5{,}000 images) as both training and
evaluation pool. The 118k-image COCO train2017 set is not currently available
on our scratch storage; rather than downloading it for this pilot we split
val2017 into two \emph{disjoint} subsets with a fixed seed
($\textsc{seed}=42$): $\text{TRAIN\_N}=4000$ images for fine-tuning and the
remaining held-out subset for evaluation
($\text{EVAL\_N}=200$ for captioning, $\text{GEN\_EVAL\_N}=30$ for generation,
the latter capped by the autoregressive cost of $576$ tokens per image). For
this study --- whose object is the \emph{relative} change of a probe under
adaptation --- this is sufficient: the same held-out subset is used to measure
the model before and after fine-tuning, so a possible mismatch in absolute
numbers vs the official benchmark splits does not affect the deltas. We will
move to the full train2017 split when scaling to the 7B model and to
benchmark-comparable numbers.

\paragraph{Captioning labels.} For an image $I$ with reference captions
$\{c_1, \dots, c_k\}$, we use the first reference $c_1$ as the training target.
Prompt and answer are masked separately and only the assistant-generated
caption tokens contribute to the loss.

\paragraph{Generation labels.} The target image is encoded to discrete
codebook indices by the model's own VQ-VAE via
\texttt{model.model.vqmodel.encode(pixel\_values).image\_tokens}; these indices
become the labels. The corresponding text prompt is the same reference caption
$c_1$. We train the generation head with teacher forcing (see
Section~\ref{sec:dryrun}).

\paragraph{LoRA configuration.}
\begin{itemize}
    \item rank $r=8$, scaling $\alpha=16$, dropout $0.05$;
    \item target modules: attention projections \texttt{q\_proj}, \texttt{k\_proj},
          \texttt{v\_proj}, \texttt{o\_proj};
    \item task type: \texttt{CAUSAL\_LM}; bias disabled;
    \item trainable parameters: $4.325$M out of $2{,}081$M total ($0.208\%$),
          distributed over $336$ LoRA tensors.
\end{itemize}

\paragraph{Optimisation.} AdamW with learning rate $10^{-4}$, batch size $1$
with gradient accumulation $8$ (effective batch size $8$), gradient clipping
to norm $1.0$. We deliberately keep batch size $1$ for this pilot: it reuses
the single-sample label masking validated by the dry-run (Section~\ref{sec:dryrun})
and side-steps any padding-related bug in batched sequences, of the kind that
caused a measurement failure in Part~I. True batching is an optimisation we
add later, after correctness.

\paragraph{Metrics.} For captioning we report CIDEr and BLEU-4 via
\texttt{pycocoevalcap} with a simple lowercase/alphanumeric tokenizer
(internally consistent for before/after deltas; PTBTokenizer will be added
for benchmark-comparable final numbers). For generation we report CLIPScore:
$100\cdot\cos(\phi_\text{img}, \phi_\text{txt})$ with
\texttt{openai/clip-vit-base-patch32}. CLIPScore is a coarse alignment metric
and may underestimate fidelity/quality loss; FID and GenEval (the two
benchmarks reported by Janus-Pro itself) are planned for the final numbers.

\paragraph{Compute.} Single A100 (40\,GB) on Snellius; \texttt{bfloat16}
weights; HF cache shared on scratch. Training one epoch on 4{,}000 examples
takes about $15$--$20$ minutes for captioning and longer for generation
(autoregressive image decoding).

\paragraph{Tracking and reproducibility.} Weights \& Biases project
\texttt{janus-peft}: training loss, learning rate, gradient norm per step,
and held-out metrics per epoch. Code, environments and SLURM scripts are
version-controlled at \texttt{github.com/kaj04/multimodal-peft}.

%% -------------------------------------------------------------------
\section{Pipeline Validation}
\label{sec:pipeline_validation}
%% -------------------------------------------------------------------

A central lesson from Part~I is that any conclusion downstream of a broken
metric is unreliable. Before any LoRA training is attempted we therefore
validate the machinery in three escalating steps: (i) inference works in both
directions, (ii) a single training step completes for every direction we want
to study, with gradients actually reaching the LoRA parameters, and (iii) the
evaluation metrics produce meaningful numbers on the untrained model. Only
after all three pass do we trust any post-training delta.

\subsection{Inference smoke test}
\label{sec:smoke}

The first check is operational: Janus-Pro-1B loads in the
\texttt{JanusForConditionalGeneration} class, captions four held-out COCO
images at a level we judge accurate by eye, and generates coherent
$384\times384$ images from text prompts. While trivial-sounding, this step
ruled out a Janus-Pro-specific failure mode we encountered along the way:
the default \texttt{generate} configuration uses \texttt{max\_length}$=20$,
which truncates the image-token sequence to far fewer than the $576$ tokens
needed for a complete picture, producing all-black images. The fix
(\texttt{max\_new\_tokens}$=576$) and the value-range handling of
\texttt{decode\_image\_tokens} (output in $[-1,1]$, HWC layout) are recorded
in Section~\ref{sec:eng_notes}.

\subsection{PEFT dry-run across all training directions}
\label{sec:dryrun}

The dry-run is the de-risking step that Anole could not pass: inject LoRA in
the live \texttt{nn.Module}, run a single training step
(forward$\rightarrow$loss$\rightarrow$backward$\rightarrow$\texttt{optimizer.step}),
and verify that gradients reach the LoRA parameters. We test each direction
in isolation, so one failure does not hide another:

\begin{enumerate}
    \item \textbf{Understanding} (image+text$\rightarrow$text): the native
          \texttt{forward(\dots, labels)} path, with labels masked to
          $-100$ before the assistant's answer and at image-token positions.
    \item \textbf{Text-only} (text$\rightarrow$text): the same native path
          without an image.
    \item \textbf{Generation} (text$\rightarrow$image): \emph{not} supported
          by the stock \texttt{forward}, which only computes text loss. We
          implement the custom training step that mirrors the model's own
          generation path:
\end{enumerate}

\noindent For generation, the loss for a single example is computed as

\begin{enumerate}
    \item encode the target image to discrete tokens
          $z = \texttt{vqmodel.encode}(I)\in\{0,\dots,16{,}383\}^{576}$;
    \item embed the text prompt $t$ with the input embedding table, and embed
          the image tokens with
          \texttt{prepare\_embeddings\_for\_image\_generation};
    \item run the shared transformer on the concatenated, teacher-forced
          sequence $[t \,|\, z_{:-1}]$;
    \item project the hidden states corresponding to the image positions
          through \texttt{generation\_head} to obtain logits over the codebook;
    \item compute $\mathcal{L}_{\text{gen}} = \mathrm{CE}(\text{logits}, z)$.
\end{enumerate}

\begin{table}[h]
\centering
\small
\begin{tabular}{lccccl}
\toprule
Direction & loss & grad-norm (LoRA) & \#LoRA params w/ grad & status \\
\midrule
Understanding (img$+$txt$\rightarrow$txt)        & 9.40 & 4.71 & 336 / 336 & PASS \\
Text-only (txt$\rightarrow$txt)                  & 3.04 & 0.58 & 192 / 336 & PASS \\
Generation (txt$\rightarrow$img, custom loss)    & 7.98 & 0.38 & 192 / 336 & PASS \\
\bottomrule
\end{tabular}
\caption{PEFT dry-run on Janus-Pro-1B (LoRA $r{=}8$ on
\texttt{q/k/v/o\_proj}, batch size $1$, AdamW). All three training directions
produce a finite loss with non-zero gradient norm at the LoRA parameters; the
custom generation loss is therefore confirmed to be differentiable end-to-end.}
\label{tab:dryrun}
\end{table}

\paragraph{Where LoRA actually lands.} With the chosen targets the $336$
injected LoRA tensors split into $192$ inside the shared language-model
backbone and $144$ inside the SigLIP understanding encoder. The generation
path does not traverse SigLIP, so only the $192$ language-model tensors
receive a gradient when training the generation direction (cf.\
Table~\ref{tab:dryrun}). This structural fact is a lever for the
LoRA-placement question (Section~\ref{sec:janus_rqs}, RQ2): the decoupled
architecture lets us, in principle, restrict adaptation to the understanding
pathway, the generation pathway, or the shared backbone, and observe where
forgetting concentrates.

\subsection{Baseline metric validation}
\label{sec:baseline}

The final check is whether our evaluation metrics behave sensibly on the
\emph{untrained} model. We caption $100$ held-out COCO images and generate
$30$ images from real COCO captions, then score them with our pipeline.

\begin{table}[h]
\centering
\small
\begin{tabular}{lcc}
\toprule
Metric & Score & Notes \\
\midrule
CIDEr     & $0.296$        & \texttt{pycocoevalcap}, simple tokenizer \\
BLEU-4    & $0.140$        & idem \\
CLIPScore & $30.6 \pm 2.9$ & CLIP ViT-B/32, $n=30$ \\
\bottomrule
\end{tabular}
\caption{Baseline evaluation of pretrained Janus-Pro-1B (no LoRA) on a held-out
COCO subset. Numbers are sensible and serve as the reference point for all
subsequent before/after deltas.}
\label{tab:baseline}
\end{table}

The numbers are in the expected range for a 1B zero-shot model. Two honest
caveats matter for interpretation later: (i) CIDEr is depressed because the
pretrained model writes long, descriptive captions whereas COCO references
are short --- low $n$-gram overlap by stylistic mismatch, not by lack of
understanding; (ii) CLIPScore measures text$\leftrightarrow$image alignment,
not generated-image fidelity or aesthetic quality, and may therefore miss a
quality regression that is invisible to alignment. Both points are addressed
in Section~\ref{sec:janus_rqs}.

%% -------------------------------------------------------------------
\section{Preliminary Results: Cross-Modal Forgetting under LoRA}
\label{sec:prelim_forgetting}
%% -------------------------------------------------------------------

\subsection{Design: a $2\times2$ cross-modal transfer matrix}

The unified-model thesis is that one shared backbone supports both
understanding and generation. The simplest experiment that puts that thesis
under pressure with PEFT is: \emph{fine-tune the LoRA adapter on one direction
and measure the probe on both}. We organise this as a $2\times2$ matrix whose
rows index the direction we train on and whose columns index the direction we
evaluate on (Table~\ref{tab:transfer_matrix}). The diagonal entries quantify
task acquisition; the off-diagonal entries quantify cross-modal forgetting.
For each cell we report the change of the same metric on the same held-out
subset relative to the pretrained baseline (Table~\ref{tab:baseline}).

\paragraph{Adapter-off vs adapter-on as a controlled comparison.}
We deliberately measure both ``before'' and ``after'' on the same loaded
model: PEFT's \texttt{disable\_adapter()} context manager turns the LoRA
update off, leaving the base weights in place; outside that context, the
update is active. The base weights, the input pipeline, the held-out subset
and the sampling seeds are therefore identical between conditions, and the
only variable is whether the rank-$r$ correction is added. This minimises
nuisance variance and isolates the effect of the adapter.

\subsection{First row: trained on captioning}
\label{sec:row_captioning}

We trained one epoch of captioning ($\textsc{train\_n}=4{,}000$,
$\textsc{eval\_n}=200$, $\textsc{gen\_eval\_n}=30$) with the configuration
of Section~\ref{sec:janus_setup}, saved the adapter, and evaluated both
metrics on the same held-out split with the adapter disabled and enabled.

\begin{table}[h]
\centering
\small
\begin{tabular}{lccc}
\toprule
Metric & Adapter off & Adapter on & $\Delta$ \\
\midrule
Captioning CIDEr   & $0.224$ & $1.243$ & $\mathbf{+1.019}$ \\
Captioning BLEU-4  & $0.127$ & $0.365$ & $+0.238$ \\
Generation CLIPScore & $30.75$ & $31.18$ & $+0.44$ \\
\bottomrule
\end{tabular}
\caption{First row of the cross-modal transfer matrix: LoRA $r{=}8$ trained
on captioning, evaluated on captioning (target) and generation (cross-modal
probe). Same held-out images and seeds with and without the adapter. The
captioning gain is large; the generation CLIPScore moves within the baseline
standard deviation ($\pm 2.9$), i.e.\ within noise.}
\label{tab:row_captioning}
\end{table}

\paragraph{Reading the captioning gain honestly.} CIDEr moves from $0.224$ to
$1.243$ ($5.5\times$), but inspection of the generated captions before and
after shows that a large part of this gain is \emph{stylistic} adaptation
rather than improved understanding. The pretrained model produces long,
descriptive captions (``\emph{The image shows a modern kitchen with a large
island featuring a sink, a faucet, and wooden cabinetry, with a view of a
dining area in the background.}'') whereas COCO references are short
(``\emph{A woman stands in the dining area.}''). After fine-tuning the model
adopts the COCO style (``\emph{A kitchen with a sink and a table.}''), which
greatly increases $n$-gram overlap with the references. The arguably more
informative content was already there before fine-tuning. We will return to
this point when designing the final evaluation (PTBTokenizer for comparability,
human-style or VQA-style probes for content fidelity).

\paragraph{Reading the generation result.} CLIPScore changes by $+0.44$,
well within the baseline standard deviation of $\pm 2.9$ over $30$ images,
so we treat the change as \emph{not significant}: the pilot shows
\textbf{no measurable cross-modal forgetting from captioning to generation}
under LoRA-on-attention. Eyeballing the before/after generated images for
five fixed prompts supports this reading: image quality and prompt fidelity
appear preserved, and the after-images are subjectively slightly more
realistic. This is the qualitative direction predicted by the literature on
LoRA as a less-forgetting alternative to full fine-tuning; the contrast with
the strong forgetting reported by Anole under full fine-tuning is the most
interesting comparison this matrix begins to draw.

\subsection{Second row: trained on generation}
\label{sec:row_generation}

We then trained the symmetric run on the generation direction with the
custom loss of Section~\ref{sec:dryrun} (same hyperparameters: $4{,}000$
training images, $1$ epoch, $r{=}8$, batch size $1$ with gradient
accumulation), saved the adapter, and ran the same before/after probe.

\begin{table}[h]
\centering
\small
\begin{tabular}{lccc}
\toprule
Metric & Adapter off & Adapter on & $\Delta$ \\
\midrule
Generation CLIPScore & $30.75$ & $28.79$ & $\mathbf{-1.96}$ \\
Captioning CIDEr     & $0.224$ & $0.250$ & $+0.026$ \\
Captioning BLEU-4    & $0.127$ & $0.126$ & $-0.001$ \\
\bottomrule
\end{tabular}
\caption{Second row of the cross-modal transfer matrix: LoRA $r{=}8$ trained
on generation, evaluated on generation (target) and captioning (cross-modal
probe). The target CLIPScore moves $-1.96$, which is within the baseline
standard deviation ($\pm 2.9$ over $n{=}30$); the captioning probe is
essentially unchanged.}
\label{tab:row_generation}
\end{table}

\paragraph{Reading the captioning probe.} The captioning metrics move by less
than the n-gram tokenisation granularity ($\Delta$ BLEU-4 $=-0.001$,
$\Delta$ CIDEr $=+0.026$): consistent with no cross-modal forgetting from
generation to captioning. Inspection of the captions before and after
confirms this: the model still produces its baseline long-form descriptions
on the same held-out images, with cosmetic wording changes
(``\emph{The image shows\dots}'' becomes ``\emph{This image shows\dots}''),
not stylistic drift.

\paragraph{Reading the generation target.} The CLIPScore decreased by
$1.96$ points, within noise but in the unexpected direction for a target
task. This is the more nuanced finding of the pilot: with $4{,}000$ images,
one epoch, $r{=}8$ and our learning rate, the custom generation loss did not
produce a positive movement on the alignment probe. Several non-exclusive
explanations apply: cross-entropy over a $16{,}384$-way codebook is a much
sparser signal than text cross-entropy and likely needs more data and/or
more epochs; CLIPScore captures only text-image alignment and may be
insensitive to a subtler quality change; and the literature on token-based
generation training (notably the Anole blog) reports a tendency of plain
cross-entropy on image tokens to drift toward mode collapse, which would
not be picked up by an alignment metric and would require a reconstruction
or perceptual term to control. Tightening this row is a clear item for the
next phase (RQ3 rank sweep on the generation direction with stronger
metrics, Section~\ref{sec:janus_rqs}).

\subsection{Pilot transfer matrix --- summary}

\begin{table}[h]
\centering
\small
\begin{tabular}{l|cc}
\toprule
\diagbox[width=12em]{train on}{evaluate on} & Captioning (CIDEr) & Generation (CLIPScore) \\
\midrule
\textbf{Captioning} & $\Delta+1.019$ (target $\uparrow$) & $\Delta+0.44$ (within noise) \\
\textbf{Generation} & $\Delta+0.026$ (within noise)      & $\Delta-1.96$ (target, within noise) \\
\bottomrule
\end{tabular}
\caption{Pilot $2\times2$ cross-modal transfer matrix on Janus-Pro-1B with
LoRA $r{=}8$ on attention (single seed, $n_{\text{cap}}{=}200$,
$n_{\text{gen}}{=}30$). The diagonal is the trained target; off-diagonal is
the cross-modal probe (negative = forgetting). Both off-diagonal entries
sit within the baseline noise of the corresponding metric, so the pilot
shows \emph{no measurable cross-modal forgetting under LoRA on attention};
the captioning target is strongly acquired (in large part as a stylistic
adaptation to COCO references), while the generation target does not move
in the expected direction with the current hyperparameters.}
\label{tab:transfer_matrix}
\end{table}

% Note for typeset: \diagbox requires \usepackage{diagbox}; if not in Ivan's
% preamble, replace the diagbox cell with a two-line header.

\paragraph{Honest scope of these numbers.} This is a single pilot run at the
1B scale with one seed and small evaluation sets ($n{=}200$ for captioning,
$n{=}30$ for generation). It is a working signal, not a definitive result.
The most informative finding is qualitative: LoRA on attention behaves as
an insulator in both directions of the unified model, in contrast with the
strong forgetting reported by Anole under full fine-tuning. The least
informative is the generation diagonal: the target metric did not move in
the expected direction, indicating that the current configuration is not
sufficient to drive useful adaptation in that direction --- a focus point
for the next experiments (Section~\ref{sec:janus_rqs}).

%% -------------------------------------------------------------------
\section{Updated Research Questions and Directions}
\label{sec:janus_rqs}
%% -------------------------------------------------------------------

The Part~I framing (Section~1) listed three RQs on dense-prediction
modalities. The Janus pivot recasts them around the
understanding$\leftrightarrow$generation axis and adds a question that only
a decoupled-encoder architecture lets us ask. The three lines below are what
the remaining experiments target. The broader threads in
Section~\ref{sec:related_work} --- token-modality-conditioned LoRA and joint
scaling laws $L_{\text{adapt}}(N, r, D_{\text{ft}}, M)$ for native
multimodal models --- remain motivating future directions; this work
is the empirical groundwork they would build on.

\paragraph{RQ1. Cross-modal forgetting under PEFT.}
\emph{Does LoRA fine-tuning on one direction (understanding or generation)
degrade the other?} Operationalised as the $2\times2$ transfer matrix of
Section~\ref{sec:prelim_forgetting}, with the first row already populated
and the second row in progress. The pilot suggests LoRA on attention behaves
as an ``insulator'' in the captioning$\rightarrow$generation direction --- a
claim we tighten with multiple seeds and stronger metrics.

\paragraph{RQ2. Localisation of forgetting along Janus' decoupled pathways.}
\emph{Is forgetting concentrated in the shared backbone, in the understanding
encoder (SigLIP), or in the generation pathway (VQ-VAE side)?} Section~\ref{sec:dryrun}
already showed that with \texttt{q/k/v/o} targets LoRA naturally lands on
both the language model ($192$ tensors) and SigLIP ($144$); the generation
forward only engages the LM half. This gives three placement regimes to
compare: LoRA restricted to the LM only, to SigLIP only, or to both
(default). RQ2 measures how each placement trades off task acquisition
against forgetting on the other direction. This is the question that the
Janus architecture uniquely makes accessible among current open native
multimodal models.

\paragraph{RQ3. Rank vs efficiency vs forgetting.}
\emph{How does LoRA rank shape the trade-off?} A sweep
$r\in\{4, 8, 16, 32, 64\}$ on the captioning direction with the same probe
suite gives a curve in (trainable parameters, target metric, cross-modal
forgetting). The shape of the curve is itself the result: a saturating gain
with growing forgetting would corroborate the ``LoRA past a small rank gives
diminishing returns'' regime reported by \citet{zhang2024scalingmeetsfinetuning};
a non-monotone forgetting curve would isolate a sweet spot.

\paragraph{Methodological strengthening planned for all three RQs.}
\begin{itemize}
    \item Multiple seeds ($\geq 3$) per cell, with error bars on every delta.
    \item A larger generation probe set than the current $n=30$.
    \item FID and GenEval to complement CLIPScore on the generation side
          (these are the benchmarks reported by Janus-Pro itself, so the
          numbers become directly comparable).
    \item PTBTokenizer for CIDEr/BLEU to be comparable with standard COCO
          captioning leaderboards.
    \item A text-perplexity probe (e.g.\ short held-out captions or
          \texttt{WikiText}) as a third forgetting axis, since both
          directions touch the language model.
\end{itemize}

\paragraph{Scaling, last.} The 7B variant of Janus-Pro and the full COCO
train2017 split are reserved for the final numbers, once the winning
configurations from RQ2 and RQ3 are identified on the 1B pilot.
The current programme is intentionally inverted with respect to scale:
finish the science on the small model, then scale only the configurations
that matter.

%% -------------------------------------------------------------------
\section{Engineering Notes and Reproducibility}
\label{sec:eng_notes}
%% -------------------------------------------------------------------

This section is the project's running list of non-obvious issues and fixes
we have hit while building the pipeline. They are tedious to rediscover, so
we keep them in writing.

\paragraph{Separate conda environments per model family.} The three model
families we touched pin incompatible \texttt{transformers} versions:
4M's \texttt{ml-4m} stack, Anole's custom fork (which patches
\texttt{transformers} 4.42.0.dev0), and Janus (recent official
\texttt{transformers}, requiring \texttt{torch} $\geq 2.5$ for
\texttt{torch.distributed.tensor}). Installing them in a single environment
breaks at least one of them. We keep three conda envs (\texttt{4m-peft},
\texttt{anole}, \texttt{janus}) and a small \texttt{env.sh} per project to
load modules and activate the right one.

\paragraph{Janus generation needs an explicit \texttt{max\_new\_tokens}.}
The default \texttt{generation\_config.max\_length}$=20$ silently truncates
the image-token sequence to far fewer than the $576$ tokens needed for a
complete image; the VQ-VAE then decodes the truncated sequence to an
all-black image, with no warning visible at the call site. We set
\texttt{max\_new\_tokens=576} for any text-to-image generation.

\paragraph{Decoding \texttt{decode\_image\_tokens} output.}
\texttt{model.decode\_image\_tokens(\dots)} returns a tensor in $[-1, 1]$,
HWC layout, that requires \texttt{grad} and lives on CUDA. Converting it to
a PIL image safely requires: (i) \texttt{detach().cpu()} before any NumPy
conversion (otherwise the conversion in \texttt{processor.postprocess} fails
because of grad/device); (ii) a single, ordered denormalisation step --- map
$[-1,1]\to[0,1]$, then clip, then $\times 255 \to$ \texttt{uint8} --- written
in one place. We centralise this in
\texttt{eval/metrics.py::decoded\_to\_pil} and route all image-saving
through it. Splitting the logic across two \texttt{if} branches with
inconsistent thresholds produced all-black outputs early in the project; the
single-path version eliminates that class of bug.

\paragraph{Batch size $1$ + gradient accumulation for the pilot.} We keep
\texttt{batch\_size}$=1$ and use gradient accumulation to reach an effective
batch size of $8$. This deliberately reuses the single-sample label
masking validated in the dry-run, and side-steps the kind of cross-sample
padding bug that, in Part~I, briefly produced an inconsistent train/val
loss before we caught it. Batched training is a future optimisation.

\paragraph{Output paths visible in Open OnDemand.} On Snellius, the
OnDemand file browser exposes the home directory but not
\texttt{/scratch-shared}. We therefore write inspectable artefacts (sample
captions, generated images, JSON results) under
\texttt{\textasciitilde/multimodal-peft/janus/.../output/} while leaving
large data (model cache, datasets) on scratch.

\paragraph{Compute-node internet.} The Snellius GPU nodes do reach the
public internet, so HuggingFace model downloads and W\&B online logging
work directly from inside a job. No offline-then-sync workflow is needed,
contrary to what we had to do on the 4M phase.

\paragraph{Weights \& Biases.} \texttt{wandb login} stores credentials in
\texttt{\textasciitilde/.netrc}, which is shared across conda envs;
re-using the existing 4M-phase credentials in the new \texttt{janus} env
required only \texttt{pip install wandb} (with the env's \texttt{python -m
pip} to avoid the system-Python install path).

\paragraph{Repository layout.} The repository
\texttt{github.com/kaj04/multimodal-peft} mirrors the project structure:
\texttt{ml-4m/} (submodule), \texttt{chameleon/} and \texttt{anole/} for
the archived exploration code, \texttt{janus/\{tests,eval,train,data\_prep\}}
for the active line, \texttt{report/} for this document. Each model family
has a \texttt{tests/} folder with smoke-test scripts mirrored by a SLURM
launcher with the same name. The single source of truth for image decoding
and metric computation lives in \texttt{janus/eval/metrics.py}.
